%% ============================================================================
%% Fichier  : annexes/Ann_B_outils.tex
%% Titre    : Référentiel des outils
%% Version  : 1.0 -- novembre 2026
%% ============================================================================

\chapter{Référentiel des Outils}
\label{ann:outils}

Les outils cités sont des exemples. Le cours insiste sur la méthode et la
compréhension~: l'outil sert le raisonnement, il ne le remplace pas. Tout
usage se fait dans un cadre autorisé ou en laboratoire fermé
(Ch.~\ref{chap:labo}).

\begin{tableaucours}{L{4.2cm}L{5.4cm}L{2.8cm}}{%
    \tableauHeader{Catégorie d'outil} &
    \tableauHeader{Usage} &
    \tableauHeader{Chapitre, lab}
}
Sources ouvertes (OSINT) & Journaux de certificats, moteurs de recherche, registres publics & Ch.~\ref{chap:osint}, HE2-LAB-02 \\
Balayage réseau & Découverte des hôtes et des services & Ch.~\ref{chap:balayage}, HE2-LAB-03 \\
Scanner de vulnérabilités & Détection assistée, à vérifier & Ch.~\ref{chap:vulnerabilites}, HE2-LAB-04 \\
Relais web local & Observation et modification des requêtes & Ch.~\ref{chap:web}, HE2-LAB-08 \\
Analyse d'annuaire & Groupes, comptes, chemins d'attaque & Ch.~\ref{chap:ad}, HE2-LAB-05 \\
Audit de mots de passe & Cassage hors ligne autorisé & Ch.~\ref{chap:mdp}, HE2-LAB-07 \\
Référentiel ATT\&CK & Planifier et mesurer la détection & Ch.~\ref{chap:detection}, HE2-LAB-09 \\
Tableur & Priorisation, calcul de gravité, matrices & Ch.~\ref{chap:vulnerabilites}, HE2-LAB-10 \\
Laboratoire virtuel & Entraînement en environnement fermé & Ch.~\ref{chap:labo} \\
\end{tableaucours}

\begin{boxalert}
\textbf{Rappel.} Aucun de ces outils ne se dirige vers un système réel sans
autorisation écrite. En dehors d'une mission mandatée, l'entraînement se fait
en laboratoire isolé (Ch.~\ref{chap:labo}).
\end{boxalert}
